\chapter{数学结构与对称变换}
\label{chap:math-structures}

同一个矩阵既可以是线性映射的坐标表示，也可以是某个变换群的元素。
如果没有先说明它作用于哪个集合、允许怎样复合、哪些性质在变换下不变，
“矩阵可逆”或“两个变换相同”就可能指完全不同的断言。
先把集合与映射说清楚，才能安全地进入群、表示和自旋量。

\begin{definition}[映射、单射、满射与双射]
给定集合 \(X,Y\)，映射 \(f:X\to Y\) 要求每个 \(x\in X\) 对应唯一的 \(f(x)\in Y\)。若 \(f(x_1)=f(x_2)\) 蕴含 \(x_1=x_2\)，称 \(f\) 为单射；若每个 \(y\in Y\) 都是某个 \(f(x)\)，称为满射；两者兼有称为双射。
\end{definition}

这一定义看似基础，却会立即工作。写逆算子 \(A^{-1}\) 时，必须先说明 \(A\) 在所讨论的定义域和值域之间为双射；只验证齐次方程没有非零解，只证明了单射。

\begin{definition}[等价关系与商集]
集合 \(X\) 上的关系 \(\sim\) 若满足自反、对称和传递，就称为等价关系。点 \(x\) 的等价类为
\[
[x]=\{y\in X:y\sim x\},
\]
所有等价类的集合记为 \(X/{\sim}\)。
\end{definition}

例如实数角度满足 \(\theta\sim\phi\) 当且仅当
\(\theta-\phi=2\pi n\)（\(n\in\mathbb Z\)）；
\([\theta]\) 表示同一个平面方向的所有角度写法。
这里有一项以后反复用到的事实：两个等价类要么完全相同，要么没有公共元素。
事实上，若 \(z\in[x]\cap[y]\)，则 \(x\sim z\) 且 \(z\sim y\)，
传递性给出 \(x\sim y\)；再对任意 \(u\sim x\) 使用传递性，就有
\(u\sim y\)。反向同理。因此商集确实是把原集合划分为互不重叠的块，
而不是随意挑出的若干代表点。后面写 \(G/N\) 时，群元会按这种方式成块。

\begin{definition}[向量空间与线性映射]
这里标量域取 \(\mathbb K=\mathbb R\) 或 \(\mathbb C\)。其上的向量空间
\(V\) 带有加法和数乘：\((V,+)\) 是交换群，并且对
\(u,v\in V\)、\(a,b\in\mathbb K\) 有
\[
a(u+v)=au+av,\quad (a+b)u=au+bu,\quad
a(bu)=(ab)u,\quad 1u=u.
\]
这四式保证“先组合向量”与“先计算系数”不发生冲突。
映射 \(T:V\to W\) 若满足
\[
T(\alpha u+\beta v)=\alpha T(u)+\beta T(v)
\]
就称为线性映射。
\end{definition}

例如 \(T:\mathbb R^2\to\mathbb R^2\)、\(T(x,y)=(x+y,y)\) 是线性映射，
但 \((x,y)\mapsto(x+y+1,y)\) 不是，因为线性映射必须把零向量送到零向量。
定义 \(\ker T=\{v\in V:Tv=0\}\) 后，便能用一个方程判定映射有没有
把两个不同输入压到同一输出。

\begin{proposition}[核与单射]
线性映射 \(T:V\to W\) 为单射，当且仅当 \(\ker T=\{0\}\)。
\end{proposition}
\begin{proof}
若 \(T\) 单射且 \(Tv=0=T0\)，则 \(v=0\)。反之，若 \(Tu=Tv\)，线性性给
\(T(u-v)=0\)，核平凡便推出 \(u=v\)。
\end{proof}

群只要求能够连续做变换和撤销变换，不要求元素可以相加；
代数则先是向量空间，再有与加法相容的乘法。
这两个概念分别把“连续做两次变换”和“生成元的乘积”说清楚。

\section{从可逆操作到群}
\label{sec:math-group-from-operations}

在力学中，绕固定轴转过角度 \(\theta\) 后再转过 \(\phi\)，等于一次转过
\(\theta+\phi\)。连续操作可以复合，不操作也算一种操作，每次转动可以撤销。
群把这些事实和复合的结合律写成公理。

\begin{definition}[群]
群是非空集合 \(G\) 连同二元运算
\(G\times G\to G,(g,h)\mapsto gh\)，使对任意 \(g,h,k\in G\)，
\((gh)k=g(hk)\)；存在 \(e\in G\) 满足 \(eg=ge=g\)；每个 \(g\) 有
\(g^{-1}\in G\) 满足 \(gg^{-1}=g^{-1}g=e\)。若 \(gh=hg\)，称为交换群。
\end{definition}

整数在加法下是群，单位元为 \(0\)，逆元为相反数。非零实数在乘法下也是群，
但全体实数在乘法下不是群，因为 \(0\) 没有逆元。
\(n\) 阶可逆矩阵在矩阵乘法下构成
\(\operatorname{GL}(n,\mathbb R)\)；矩阵乘法一般不交换。
例如取
\(A=\begin{psmallmatrix}1&1\\0&1\end{psmallmatrix}\)、
\(B=\begin{psmallmatrix}1&0\\1&1\end{psmallmatrix}\)，则
\(AB=\begin{psmallmatrix}2&1\\1&1\end{psmallmatrix}\)
而 \(BA=\begin{psmallmatrix}1&1\\1&2\end{psmallmatrix}\)。
这不是记号问题，而是两次操作的顺序真的会改变结果。

\begin{proposition}[单位元、逆元与消去律]
群的单位元唯一，每个群元的逆元唯一。若 \(gh=gk\)，则 \(h=k\)；
并且 \((gh)^{-1}=h^{-1}g^{-1}\)。
\end{proposition}
\begin{proof}
若 \(e,e'\) 都是单位元，则 \(e=ee'=e'\)。若 \(u,v\) 都是 \(g\) 的逆元，
结合律给出 \(u=u(gv)=(ug)v=v\)。由 \(gh=gk\) 左乘 \(g^{-1}\) 得
\(h=k\)。最后直接计算
\((gh)(h^{-1}g^{-1})=e=(h^{-1}g^{-1})(gh)\)，再用逆元唯一性。
\end{proof}

讨论较小的一组允许操作时，还需要知道它本身是否仍是群。
\begin{definition}[子群]
若非空子集 \(H\subseteq G\) 对乘法及取逆封闭，则 \(H\) 连同原运算
称为 \(G\) 的子群，记作 \(H\leq G\)。
\end{definition}
例如绕同一轴转动所得的矩阵是三维转动群的子群；在正实数乘法群中，
只取大于 \(1\) 的伸缩因子则不是子群，因为逆因子小于 \(1\)。
定义中的单位元不必另加：取 \(h\in H\)，就有
\(hh^{-1}=e\in H\)。

一个有限群里不可能把子群任意大小地塞进去。给定 \(g\in G\)，把
\(H\) 的每个元素都左乘 \(g\)，得到集合
\(gH=\{gh:h\in H\}\)，称为一个左陪集。左乘 \(g\) 可由左乘
\(g^{-1}\) 撤销，所以 \(gH\) 与 \(H\) 元素个数相同。
更重要的是，两个左陪集只要有一个公共元素就必须相同：若
\(gh_1=kh_2\)，则 \(k^{-1}g=h_2h_1^{-1}\in H\)，从而任意
\(h\in H\) 都满足 \(gh=k(h_2h_1^{-1}h)\in kH\)。
交换 \(g,k\) 再做一次就得到 \(gH=kH\)。

\begin{theorem}[Lagrange 定理]
若有限群 \(G\) 的子群为 \(H\)，则 \(|H|\) 整除 \(|G|\)，且
\[
|G|=[G:H]|H|,
\]
其中 \([G:H]\) 表示不同左陪集的个数，称为 \(H\) 在 \(G\) 中的指数。
\end{theorem}
\begin{proof}
每个 \(g\) 属于 \(gH\)，所以全部左陪集覆盖 \(G\)；上一段证明
它们互不相交或完全相同，而且每块都有 \(|H|\) 个元素。
逐块计数便得等式。
\end{proof}

若 \(g\in G\)，反复做它所成的集合
\(\langle g\rangle=\{g^k:k\in\mathbb Z\}\) 是子群，称为由 \(g\)
生成的循环子群。有限群中存在最小正整数 \(m\) 使 \(g^m=e\)：
因为幂 \(e,g,g^2,\ldots\) 总会重复，消去后就有一个正幂等于单位元。
此时 \(e,g,\ldots,g^{m-1}\) 两两不同，所以
\(|\langle g\rangle|=m\)，即元素 \(g\) 的阶为 \(m\)。
Lagrange 定理于是推出 \(g^{|G|}=e\)。若 \(|G|\) 是素数，任取
\(g\ne e\)，其阶既大于一又整除 \(|G|\)，只能等于 \(|G|\)；
所以任意素数阶群都是循环群。这些结论不要求乘法交换。

\section{把抽象操作作用到具体对象}

知道两个转动怎样复合，还不足以知道它们怎样改变位置、速度或量子态。
后一个问题需要指定操作所作用的集合。
\begin{definition}[群作用]
群 \(G\) 在集合 \(X\) 上的左作用是映射
\(G\times X\to X,(g,x)\mapsto g\cdot x\)，满足
\[
e\cdot x=x,\qquad (gh)\cdot x=g\cdot(h\cdot x).
\]
点 \(x\) 的轨道为 \(G\cdot x=\{g\cdot x:g\in G\}\)，稳定子为
\(G_x=\{g\in G:g\cdot x=x\}\)。
\end{definition}

第二条公理固定顺序：\(gh\) 先把 \(h\) 作用到 \(x\)，再作用 \(g\)。
以平面旋转为例，原点的轨道只有原点且稳定子是整个旋转群；非零向量的轨道
是一个圆，稳定子只有恒等旋转。稳定子是子群，因为保持 \(x\) 不动的操作
的复合与逆仍保持 \(x\) 不动。

\begin{theorem}[轨道与稳定子的计数]
若 \(G\) 是有限群并作用在 \(X\) 上，则对任意 \(x\in X\)，
\[
|G\cdot x|=\frac{|G|}{|G_x|}.
\]
\end{theorem}
\begin{proof}
固定 \(y=g\cdot x\)。另一个 \(h\in G\) 满足 \(h\cdot x=y\)，当且仅当
\(g^{-1}h\cdot x=x\)，即 \(h\in gG_x\)。因此映射
\(g\mapsto g\cdot x\) 在轨道每个点上恰有 \(|G_x|\) 个原像。
这些原像互不重叠且覆盖 \(G\)，故 \(|G|=|G\cdot x|\,|G_x|\)。
\end{proof}

这个证明也说明为什么不能凭“有 \(|G|\) 个操作”就断言产生 \(|G|\) 个
不同状态：稳定子中的操作在给定 \(x\) 上产生同一个结果。
正三角形给出一个看得见的核对。把三个顶点标为 \(1,2,3\)，其刚性
对称群记为 \(D_3\)：有恒等操作、绕中心转 \(2\pi/3\) 与
\(4\pi/3\) 的两个转动，以及三条穿过顶点和对边中点的反射。
取 \(x=1\)，六个操作把它送到的不同位置只有 \(\{1,2,3\}\)，
所以轨道有三个元素；保持顶点 \(1\) 不动的恰是恒等操作和穿过它的
反射，稳定子有两个元素。式中的 \(6=3\times2\) 不是巧合：
每一个目标顶点都被恰好两个操作送达。

群作用还让我们分清“抽象操作不同”和“观测结果不同”。
若实验只读取一个顶点的位置，稳定子中的两次操作无法由这个读数区分；
若读出整个三角形的取向，原来的反射便不再与恒等操作相同。
因此轨道与稳定子总是针对一个明确的作用集合和一个明确的点而言。

\section{表示、同态与生成元}

量子态可以线性叠加，因此对称性作用到态空间时也应与线性结构相容。
\begin{definition}[群同态与线性表示]
群同态 \(f:G\to H\) 满足 \(f(gh)=f(g)f(h)\)。其核为
\(\ker f=\{g:f(g)=e_H\}\)。若 \(V\) 是向量空间，值域为可逆线性映射群
\(\operatorname{GL}(V)\) 的同态
\(\rho:G\to\operatorname{GL}(V)\) 称为 \(G\) 在 \(V\) 上的表示。
\end{definition}

同态关系给出 \(f(e_G)=e_H\) 和 \(f(g^{-1})=f(g)^{-1}\)。
若一个操作位于 \(\ker\rho\)，它在这个表示中作用为恒等，
却不一定是抽象群的单位元。若 \(\ker\rho=\{e\}\)，表示称为忠实表示。

\begin{proposition}[像、核与单射]
同态 \(f:G\to H\) 的像 \(f(G)\) 是 \(H\) 的子群，核是 \(G\) 的正规子群；
这里正规意为对所有 \(g\in G,k\in\ker f\)，都有 \(gkg^{-1}\in\ker f\)。
并且 \(f\) 单射，当且仅当 \(\ker f=\{e_G\}\)。
\end{proposition}
\begin{proof}
像对乘法和逆元封闭。若 \(k\in\ker f\)，则
\(f(gkg^{-1})=f(g)e_Hf(g)^{-1}=e_H\)，故核正规。
若 \(f(g)=f(h)\)，则 \(f(h^{-1}g)=e_H\)；核平凡时推出 \(g=h\)。
反方向取 \(f(g)=e_H=f(e_G)\) 即可。
\end{proof}

\section{什么时候可以把不同操作视为同一个操作}

一个表示可能有非平凡的核：若 \(k\in\ker\rho\)，
\(g\) 和 \(gk\) 对所有被表示的状态作用完全相同。
要把它们当作同一个有效操作，不能只说“忽略 \(k\)”；
必须证明所得乘法仍有确定的结果。

\begin{definition}[正规子群与商群的候选元素]
子群 \(N\leq G\) 若对所有 \(g\in G\) 有 \(gNg^{-1}=N\)，
称为正规子群，记作 \(N\trianglelefteq G\)。
它与 \(gN=Ng\) 等价：把 \(gNg^{-1}=N\) 右乘 \(g\) 即得后式，
反过来再右乘 \(g^{-1}\) 也成立。
所有左陪集组成的集合记为 \(G/N\)。尝试定义两类的乘法为
\[(gN)(hN)=(gh)N.\]
\end{definition}

这里“尝试”二字不可省略。若把 \(g\) 改写成同一陪集里的 \(gn\)，
结果不能因这个写法而变；这正是正规性要解决的问题。

\begin{theorem}[商群的良定义性]
上面的乘法不依赖所选代表元，当且仅当 \(N\) 为正规子群。
此时 \(G/N\) 成群，单位元为 \(N\)，\((gN)^{-1}=g^{-1}N\)。
\end{theorem}
\begin{proof}
先设 \(N\) 正规。若 \(g'=gn_1\)、\(h'=hn_2\)，
则
\[
g'h'=gn_1hn_2
=gh\,(h^{-1}n_1h)n_2.
\]
正规性使括号中的元素属于 \(N\)，故 \(g'h'N=ghN\)。
乘法因而良定义；结合律、单位元和逆元直接由 \(G\) 的公理传下。
反过来，若陪集乘法良定义，则 \(gN=g n N\) 对所有 \(n\in N\)；
分别乘 \(g^{-1}N\) 并比较代表元，等价地从
\((gN)(nN)(g^{-1}N)=(gN)(N)(g^{-1}N)=N\)
得到 \(gng^{-1}\in N\)。所以 \(N\) 正规。
\end{proof}

\begin{theorem}[同态的第一同构定理]
对群同态 \(f:G\to H\)，令 \(N=\ker f\)。
则 \(\bar f:G/N\to f(G)\)、\(\bar f(gN)=f(g)\)
是群同构，即保持乘法的双射。
\end{theorem}
\begin{proof}
前一命题已证明 \(N\) 正规。若 \(gN=hN\)，
则 \(h^{-1}g\in N\)，从而 \(f(g)=f(h)\)，所以 \(\bar f\)
不依赖代表元。它保持乘法，因为
\(\bar f((gN)(hN))=f(gh)=f(g)f(h)\)。
到 \(f(G)\) 的满射性由像的定义给出；
若 \(\bar f(gN)=\bar f(hN)\)，则 \(h^{-1}g\in N\)，
即 \(gN=hN\)，故它也单射。
\end{proof}

例如加法群 \(\mathbb Z\) 的子群 \(n\mathbb Z\) 正规，
\(\mathbb Z/n\mathbb Z\) 的元素是整数的余数类，
而不是任选的 \(0,\ldots,n-1\) 那些代表数。
取模映射 \(k\mapsto k\bmod n\) 的核正是 \(n\mathbb Z\)；
同构定理解释了为什么余数相加能定义一个群。
下一节的 \(SU(2)\) 与三维转动也会出现完全相同的结构，
只是核从无限个整数变成 \(\{I,-I\}\)。

平面旋转的向量表示为
\[
R(\theta)=
\begin{pmatrix}\cos\theta&-\sin\theta\\
\sin\theta&\cos\theta\end{pmatrix},\qquad
R(\theta+\phi)=R(\theta)R(\phi).
\]
矩阵乘法中的加角公式验证了表示公理。取 \(\theta=0\) 处导数，得到
\[
J=\left.\frac{\dd R(\theta)}{\dd\theta}\right|_{\theta=0}
=\begin{pmatrix}0&-1\\1&0\end{pmatrix},\qquad J^2=-I.
\]
由矩阵指数的幂级数，把偶次和奇次幂分别相加，就得到
\(\ee^{\theta J}=I\cos\theta+J\sin\theta=R(\theta)\)。
连续转动由一个线性算子的指数生成。

量子态的范数要在对称变换下保持，因此矩阵表示通常取酉矩阵。
记 \(U(n)=\{M\in\mathbb C^{n\times n}:M^\dagger M=I\}\)。
\begin{definition}[有限维酉表示]
在复内积空间 \(V\) 上，群表示 \(U:G\to\operatorname{GL}(V)\)
若对每个 \(g\) 满足 \(U(g)^\dagger U(g)=I\)，
称为酉表示。
\end{definition}

\begin{proposition}[单参数酉变换的生成元]
设 \(U:\mathbb R\to U(n)\) 可微，
\(U(\theta+\phi)=U(\theta)U(\phi)\) 且 \(U(0)=I\)。
则存在唯一 Hermite 矩阵 \(G=G^\dagger\) 使
\[
U(\theta)=\exp(-\ii\theta G/\hbar).
\]
\end{proposition}
\begin{proof}
令 \(A=U'(0)\)。把群关系对 \(\phi\) 在零点求导，
得到 \(U'(\theta)=U(\theta)A\)，初值为 \(I\)；
矩阵常微分方程的幂级数解是 \(U(\theta)=e^{\theta A}\)。
再对 \(U(\theta)^\dagger U(\theta)=I\) 在零点求导，
得 \(A^\dagger+A=0\)。取 \(G=\ii\hbar A\)，
便有 \(G^\dagger=G\)、\(A=-\ii G/\hbar\)。
若另有 \(G'\) 给同一个 \(U\)，比较零点导数即得 \(G'=G\)。
\end{proof}

例如转角是无量纲参数时，\(G\) 的单位与 \(\hbar\) 相同，
正是角动量生成元的单位。这个证明只针对有限维矩阵；
波函数空间里的位置和动量生成元常是无界算子，
必须先说明定义域与自伴性，不能直接照搬矩阵幂级数。

\begin{definition}[矩阵 Lie 群及其 Lie 代数]
若矩阵群 \(G\subseteq\operatorname{GL}(n,\mathbb R)\) 在每个元素附近
都能由一个开集 \(U\subseteq\mathbb R^d\) 光滑地参数化：存在映射
\(F:U\to G\)，它一一对应于 \(G\) 的一个局部邻域、逆映射连续，
而 \(F\) 的 \(d\) 个参数导数线性无关，就称它为矩阵 Lie 群。
这里“邻域”采用矩阵元的通常距离；所需的局部光滑性只涉及有限个
实变量的多元微积分。经过单位矩阵的可微群内曲线 \(g(t)\)
满足 \(g(0)=I\)；它们在零点的速度 \(X=\dot g(0)\)
构成单位元处的切空间 \(\mathfrak g\)，配上矩阵交换子后称为 Lie 代数。
\end{definition}

在参数 \(F(u)\) 中，曲线速度是 \(DF(u_0)\dot u(0)\)，所以速度集合
是 \(DF(u_0)\) 的像，确实是一个向量空间，而不只是若干曲线的杂乱集合。
光滑曲线的坐标无关定义将在流形部分给出；此处只需对矩阵元逐项求导。
交换子 \([X,Y]=XY-YX\) 与操作顺序的关系可以直接计算。
令 \(g(s)=I+sX+O(s^2)\)、
\(h(t)=I+tY+O(t^2)\)，直接相乘得到
\[
g(s)h(t)g(s)^{-1}h(t)^{-1}
=I+st(XY-YX)+O(s^2t+st^2).
\]
所以先做两个很小的操作与交换顺序之间，最低非零差别由 \([X,Y]\) 控制。
它也确实仍属于 \(\mathfrak g\)：固定群内曲线 \(g(s)\) 时，
\(g(s)h(t)g(s)^{-1}\) 对 \(t\) 在零点求导给
\(g(s)Yg(s)^{-1}\in T_I G\)。
这个向量随 \(s\) 光滑变化，故在向量空间 \(T_I G\) 内
对 \(s\) 求导仍留在其中；在 \(s=0\) 得
\(XY-YX\)。这说明“Lie 代数”不仅有向量加法，
交换子也封闭。

把这个定义用到三维转动群会得到熟悉的叉积，而不需要预先知道抽象 Lie 理论。
设
\(SO(3)=\{R\in\mathbb R^{3\times3}:R^{\mathsf T}R=I,\det R=1\}\)。
它在单位矩阵附近确有三个光滑局部参数：对充分小的反对称矩阵 \(A\) 取
\(R=(I-A)(I+A)^{-1}\)。分母在 \(A=0\) 可逆，直接相乘验证
\(R^{\mathsf T}R=I\)；而 \(R\) 连续接在 \(I\) 上，故 \(\det R=1\)。
反过来，对靠近 \(I\) 的转动，\(I+R\) 可逆且
\(A=(I-R)(I+R)^{-1}\) 反对称。
这两个有理式互为逆映射，给出了所需的局部参数，且参数个数正是
反对称三阶矩阵的三个独立矩阵元。
若 \(R(t)\in SO(3)\)、\(R(0)=I\)，对正交关系求导得
\(X^{\mathsf T}+X=0\)，其中 \(X=\dot R(0)\)。
反过来，任意实反对称矩阵 \(X\) 都产生曲线 \(e^{tX}\)：
因为 \((e^{tX})^{\mathsf T}e^{tX}=e^{-tX}e^{tX}=I\)，
且曲线从行列式一的单位矩阵连续出发，行列式始终为一。
因此 \(SO(3)\) 的 Lie 代数\emph{恰好}是三阶实反对称矩阵空间。

每个这样的矩阵可唯一写为
\[
L(\mathbf a)=
\begin{pmatrix}
0&-a_3&a_2\\ a_3&0&-a_1\\ -a_2&a_1&0
\end{pmatrix},\qquad
L(\mathbf a)\mathbf v=\mathbf a\times\mathbf v.
\]
把它作用到任意 \(\mathbf v\)，利用三重叉积恒等式可直接核对
\[
[L(\mathbf a),L(\mathbf b)]\mathbf v
=\mathbf a\times(\mathbf b\times\mathbf v)
-\mathbf b\times(\mathbf a\times\mathbf v)
=(\mathbf a\times\mathbf b)\times\mathbf v.
\]
所以 \([L(\mathbf a),L(\mathbf b)]=L(\mathbf a\times\mathbf b)\)。
取三个坐标轴的单位向量，便是
\([L_i,L_j]=\epsilon_{ijk}L_k\)。非交换性在这里有可视的含义：
绕不同轴做两次极小转动，交换顺序后剩下沿第三个方向的转动。

\begin{exercise}
正方形的四个旋转构成 \(C_4=\{e,r,r^2,r^3\}\)，其中 \(r^4=e\)。
求一个顶点的轨道和稳定子，并核对轨道与稳定子的计数公式。
\end{exercise}

\begin{exercise}
取 \(D=\operatorname{diag}(1,-1)\)。证明 \(D^2=I\)，并计算
\(DR(\theta)D^{-1}\)。这个计算说明反射如何改变旋转方向。
\end{exercise}

\section{为什么一阶方程需要一种新的乘法}
\label{sec:math-clifford-first}

相对论色散 \(E^2=m^2c^4+c^2|\vb{p}|^2\) 是动量的二次关系。
若希望像 Schrödinger 方程那样只对时间求一次导数，又希望空间导数也只出现一次，
就得寻找矩阵 \(A_i,B\)，使 \(cA_i p_i+Bmc^2\) 的平方等于原二次式乘单位矩阵。
普通数做不到：交叉项 \(p_ip_j\) 无法消失。矩阵可以，因为乘法不必交换。

\begin{definition}[二次型与极化]
实向量空间 \(V\) 上的二次型是函数 \(q(v)=g(v,v)\)，其中 \(g\) 是对称双线性型。
反过来，\(g\) 可由
\[
g(v,w)=\frac12\{q(v+w)-q(v)-q(w)\}
\]
恢复。这里双线性指每固定一个变量，对另一个变量都是线性的。
\end{definition}

例如 Euclid 空间的 \(q(v)=\sum_i v_i^2\) 给出点积
\(g(v,w)=\sum_i v_iw_i\)。若 \(q(v)=v_0^2-v_1^2-v_2^2-v_3^2\)，
则所得双线性型有一正三负的号差，正是相对论的度规形式。

\begin{definition}[结合代数与双侧理想]
域 \(\mathbb K\) 上的结合代数是 \(\mathbb K\)-向量空间 \(A\)，
另带双线性乘法 \(A\times A\to A\)，满足
\((ab)c=a(bc)\)；若存在 \(1\) 使 \(1a=a1=a\)，称它含单位元。
子空间 \(I\subset A\) 若对任意 \(a\in A,i\in I\)
都有 \(ai,ia\in I\)，称双侧理想。
\end{definition}

双侧条件保证商空间 \(A/I\) 上的乘法
\((a+I)(b+I)=ab+I\) 与代表元无关：
把 \(a,b\) 分别换成 \(a+i,b+j\)，乘积的新旧差为
\(aj+ib+ij\in I\)。这与上一节商群的良定义性
是同一个逻辑检查，只是现在还保留加法与数乘。

\begin{definition}[Clifford 代数]
给定 \((V,q)\)，在一个含单位元 \(1\) 的结合代数中加入 \(V\) 的元素，
并要求
\[
vw+wv=2g(v,w)1\qquad(v,w\in V).
\]
由这些生成元及关系得到的代数称为 \(\operatorname{Cl}(V,q)\)。
一种具体构造是先取由 \(V\) 的有限有序字串张成的向量空间，
只要求每个字串位置对向量加法和数乘保持线性；空字串作为单位元。
字串连接给出结合乘法。再把所有
\(vv-q(v)1\) 及它们左右乘任意字串所得的表达式规定为零，
即对由这些表达式生成的双侧理想取商。
\end{definition}

取商以前，字串乘法尚不附加交换律；取商就是宣布
\(v^2=q(v)1\) 必须成立。把 \((v+w)^2=q(v+w)1\) 展开，
减去 \(v^2=q(v)1\) 与 \(w^2=q(w)1\)，恰得到定义中的反对易关系。
所以关系不是额外猜测，而是二次型的极化。

\begin{proposition}[一阶算子的平方]
若矩阵 \(\gamma^\mu\) 满足
\(\gamma^\mu\gamma^\nu+\gamma^\nu\gamma^\mu=2g^{\mu\nu}I\)，
且坐标偏导可交换，则
\[
(\gamma^\mu\partial_\mu)^2
=g^{\mu\nu}\partial_\mu\partial_\nu I.
\]
\end{proposition}
\begin{proof}
把两个相同的导数指标 \(\mu,\nu\) 交换后取平均：
\[
\gamma^\mu\gamma^\nu\partial_\mu\partial_\nu
=\frac12(\gamma^\mu\gamma^\nu+\gamma^\nu\gamma^\mu)
\partial_\mu\partial_\nu
=g^{\mu\nu}\partial_\mu\partial_\nu I.
\]
第二个等号只用了偏导交换和反对易关系；若改用不对易的协变导数，
还会留下曲率或场强项，不能照搬此证明。
\end{proof}

\section{Pauli 矩阵让关系第一次工作}

取三矩阵
\[
\sigma_1=\begin{pmatrix}0&1\\1&0\end{pmatrix},\quad
\sigma_2=\begin{pmatrix}0&-\ii\\\ii&0\end{pmatrix},\quad
\sigma_3=\begin{pmatrix}1&0\\0&-1\end{pmatrix}.
\]
直接相乘可查
\(\sigma_i\sigma_j+\sigma_j\sigma_i=2\delta_{ij}I\)。
例如 \(\sigma_1\sigma_2=\ii\sigma_3\)，而
\(\sigma_2\sigma_1=-\ii\sigma_3\)；其余两个不同指标的情形
由循环更换 \(1,2,3\) 得到，同指标则是 \(\sigma_i^2=I\)。
对 \(v=(v_1,v_2,v_3)\in\RR^3\)，记
\(\sigma(v)=v_i\sigma_i\)。由上式，
\[
\sigma(v)^2=|v|^2I,\qquad
\frac12\operatorname{tr}[\sigma(v)\sigma(w)]=v\cdot w.
\]
因此一个三维向量可用无迹厄米矩阵表达，而其长度被矩阵乘法保存。
开头的一阶能量算子还需要一个与三个空间方向都反对易的质量矩阵。
把二阶矩阵拼成四阶分块矩阵即可做到：
\[
\alpha_i=\begin{pmatrix}0&\sigma_i\\\sigma_i&0\end{pmatrix},
\qquad
\beta=\begin{pmatrix}I&0\\0&-I\end{pmatrix}.
\]
分块乘法给出
\(\alpha_i\alpha_j+\alpha_j\alpha_i=2\delta_{ij}I_4\)、
\(\alpha_i\beta+\beta\alpha_i=0\)、\(\beta^2=I_4\)。
于是对相互可交换的三个动量数 \(p_i\)，定义
\(H(\vb{p})=c\alpha_i p_i+\beta mc^2\)，展开平方时
不同动量方向的交叉项和动量--质量交叉项成对抵消，严格得到
\[
H(\vb{p})^2=
\left(c^2|\vb{p}|^2+m^2c^4\right)I_4.
\]
若将 \(p_i\) 换成作用于足够光滑波函数的
\(-\ii\hbar\partial_i\)，各偏导仍可交换，故同一代数计算成立。
若再让不同 \(p_i\) 与空间变化的电磁势耦合，动量分量一般不再交换，
便不能只留下右边的标量平方；额外项正承载自旋与场的耦合。
这里尚未依赖任何具体本征态，说明相对论色散首先是一条矩阵代数恒等式。

这里需要一种不同于群表示、但同样保结构的映射：代数同态
是保持加法、数乘、单位元和乘法的映射。
把 \(\operatorname{Cl}(\RR^3,|\cdot|^2)\) 的生成元
送到相应 Pauli 矩阵时，由于矩阵满足相同关系，
这个指定会延拓为到复矩阵代数的同态。
这里的代数和同态先按\emph{实}线性理解；复矩阵只是值域，不能因为
矩阵里有 \(\ii\) 就擅自把原来的实标量域改为复数。事实上还可以证明
这个特例的同态恰为同构，而不只说它满足相同关系。

令 \(e_1,e_2,e_3\) 为 \(\RR^3\) 的标准基。关系给出
\(e_i^2=1\)、\(e_ie_j=-e_je_i\)（\(i\ne j\)）。因此任意字串
都可先把重复的 \(e_i\) 消去，再把下标从小到大排列，故整个商代数
至多由以下八个实向量张成：
\[
1,\quad e_1,e_2,e_3,\quad e_1e_2,e_2e_3,e_3e_1,
\quad e_1e_2e_3.
\]
它们的矩阵像依次是
\(I,\sigma_1,\sigma_2,\sigma_3,
\ii\sigma_3,\ii\sigma_1,\ii\sigma_2,\ii I\)。
这些矩阵作为\emph{实}向量线性无关：任意复二阶矩阵可唯一写为
\(z_0I+z_i\sigma_i\)，将每个复系数拆成实部与虚部恰得到这八项。
于是商代数确为八维，映射既单射又满射，得到
\(\operatorname{Cl}(\RR^3,|\cdot|^2)\cong M_2(\CC)\)
（作为实代数）。这次逐项核对也提醒我们：生成元的矩阵关系给出
一个表示，但一般情况下不能只凭关系相同就宣布表示忠实。

\section{为什么旋量转一圈会变号}

上面的二分量复列向量是否与普通三维向量同样变换？用一个具体转动回答。
定义
\[
SU(2)=\{U\in\CC^{2\times2}:U^\dagger U=I,\ \det U=1\}.
\]
由于 \(U\sigma(v)U^\dagger\) 仍是无迹厄米矩阵，它等于唯一的
\(\sigma(R_Uv)\)，这定义了实线性映射 \(R_U:\RR^3\to\RR^3\)。

\begin{proposition}[从二分量矩阵到空间转动]
映射 \(U\mapsto R_U\) 是到 \(SO(3)\) 的群同态，核恰为 \(\{I,-I\}\)。
\end{proposition}
\begin{proof}
两次共轭给 \(R_{UV}=R_UR_V\)。由迹的循环性质，
\[
R_Uv\cdot R_Uw
=\frac12\operatorname{tr}[U\sigma(v)U^\dagger
U\sigma(w)U^\dagger]
=v\cdot w.
\]
故 \(R_U\) 正交。事实上，把酉矩阵第一列写成
\((z,w)^{\mathsf T}\)，由两列正交且行列式为一，第二列
只能是 \((-w^*,z^*)^{\mathsf T}\)，且 \(|z|^2+|w|^2=1\)。
分别取 \(z,w\) 的实部和虚部，便得
\(U=aI-\ii\vb{b}\cdot\vb{\sigma}\)，其中
\(a\) 为实数、\(\vb{b}\in\RR^3\)，且 \(a^2+|\vb{b}|^2=1\)；
令 \(\alpha=\arccos a\in[0,\pi]\)，在 \(\vb{b}\ne0\) 时取
\(\vb{n}=\vb{b}/|\vb{b}|\)，在 \(\vb{b}=0\) 时任选单位向量
\(\vb{n}\)。则
\(U(t)=\cos(t\alpha)I-\ii\sin(t\alpha)\,
\vb{n}\cdot\vb{\sigma}\)（\(0\leq t\leq1\)）
是一条从 \(I\) 到 \(U\) 的 \(SU(2)\) 内连续路径。
正交矩阵的行列式只能为 \(\pm1\)，故 \(\det R_U=1\)。
若 \(R_U=I\)，则 \(U\) 与三个 \(\sigma_i\) 都对易；
与 \(\sigma_3\) 对易迫使 \(U\) 为对角矩阵，再与 \(\sigma_1\) 对易
迫使两个对角元相等。于是 \(U=\lambda I\)；
酉性与 \(\det U=1\) 给 \(\lambda=\pm1\)。
\end{proof}

还要确认所有空间转动都能这样得到。实正交三阶矩阵的特征值
模长为一，非实根成共轭对；若有一对非实根，它们的积为一，
剩下的实根只能为一。若三个根都实，它们各为 \(\pm1\)，
行列式为一也迫使至少一个根为一。
其单位特征向量 \(\vb{n}\) 是不动轴。
\(\vb{n}^\perp\) 也保持不变，而限制在这个二维平面上的
矩阵行列式为一，故是转过某角 \(\theta\) 的平面转动。
取
\[
U(\theta,\vb{n})
=\cos\frac{\theta}{2}I-\ii\sin\frac{\theta}{2}\,
\vb{n}\cdot\vb{\sigma} .
\]
由 \(\sigma_i\sigma_j=\delta_{ij}I+
\ii\epsilon_{ijk}\sigma_k\) 先算出
\[
(\vb{n}\cdot\vb{\sigma})\sigma(v)
(\vb{n}\cdot\vb{\sigma})
=2(\vb{n}\cdot v)\sigma(\vb{n})-\sigma(v).
\]
把 \(U\) 的两项展开，得到
\[
U\sigma(v)U^\dagger
=\sigma\!\left(v\cos\theta+
(\vb{n}\times v)\sin\theta+
\vb{n}(\vb{n}\cdot v)(1-\cos\theta)\right).
\]
括号中正是绕 \(\vb{n}\) 转 \(\theta\) 的 Rodrigues 公式，
故每个 \(SO(3)\) 转动都有原像，映射满射。
但 \(U(2\pi,\vb{n})=-I\)，而 \(R_U=I\)。
第一同构定理现在给出精确的关系
\[
SU(2)/\{I,-I\}\cong SO(3).
\]
左边的一个元素是 \(\{U,-U\}\) 这个陪集，
而不是单独的 \(U\)。两者对三维向量作同一个转动，
但在二分量状态空间上的作用不同。
二分量列向量在 \(U\) 下变化时称为旋量：一次 \(2\pi\) 转动使旋量变号，
向量却回到原值。这个负号在单个态的概率中不可见，却可在相干比较中留下相位。

“留下相位”可以直接算，而不必把旋量变号误当作单一路径的可观测量。
设 \(|A\rangle,|B\rangle\) 是正交的两条路径态，\(|\chi\rangle\) 是归一化
二分量自旋态。入射态取
\(\bigl(|A\rangle+|B\rangle\bigr)|\chi\rangle/\sqrt2\)。
只在路径 \(B\) 上完成一次 \(2\pi\) 旋转后，末态变为
\(\bigl(|A\rangle-|B\rangle\bigr)|\chi\rangle/\sqrt2\)。
若把两路重新合并并投影到对称路径态
\(|+\rangle=(|A\rangle+|B\rangle)/\sqrt2\)，
原来这个出口的概率为 \(1\)，旋转后为 \(0\)。
因此可观测的是\emph{两条路径之间}的相对负号，而不是孤立旋量的整体相位。

最后把有限转动与上一节的生成元对应起来。若 \(U(t)\in SU(2)\)、
\(U(0)=I\)，其速度 \(X=\dot U(0)\) 满足
\(X^\dagger+X=0\)，这是对 \(U^\dagger U=I\) 求导所得；
对 \(\det U=1\) 求导又给 \(\operatorname{tr}X=0\)，因为行列式
在单位矩阵处的一阶变化只由对角矩阵元之和贡献。
反过来，每个这样的二阶矩阵唯一写成
\(X=-\ii b_i\sigma_i/2\)（\(b_i\in\RR\)）。
利用 \((\vb{b}\cdot\vb{\sigma})^2=|\vb{b}|^2I\)
按奇偶幂相加，在 \(\vb{b}\ne0\) 时 \(\ee^{tX}\) 恰是上面的
\(U(t|\vb{b}|,\vb{b}/|\vb{b}|)\)；\(\vb{b}=0\) 时它就是
恒等矩阵。因此每个这样的速度都对应 \(SU(2)\) 内的曲线。
这证明 \(SU(2)\) 的 Lie 代数正是三维实向量空间
\(\operatorname{span}_{\RR}\{-\ii\sigma_i/2\}\)。

从显式矩阵乘法还得到
\([\sigma_i,\sigma_j]=2\ii\epsilon_{ijk}\sigma_k\)。
记 \(T_i=-\ii\sigma_i/2\)，则
\([T_i,T_j]=\epsilon_{ijk}T_k\)，与空间转动的
\([L_i,L_j]=\epsilon_{ijk}L_k\) 完全一致。
若把 \(S_i=\hbar\sigma_i/2\) 作为有角动量单位的自旋算符，
同一计算变成 \([S_i,S_j]=\ii\hbar\epsilon_{ijk}S_k\)。
因此二分量旋量与三维向量在\emph{无穷小}转动上遵守同一套复合规则，
却在完整转过 \(2\pi\) 后给出不同结果；只看生成元的交换关系，
不能看出这种全局的正负号。

\begin{exercise}
取 \(\vb{n}=(0,0,1)\)，直接算
\(U(\theta,\vb{n})\sigma_1U(\theta,\vb{n})^\dagger\)，
核对得到 \(\cos\theta\,\sigma_1+\sin\theta\,\sigma_2\)。
\end{exercise}


等价类上的公式还需要额外小心：若用代表元 \(x\) 定义
\(F([x])=f(x)\)，必须证明 \(x\sim y\) 时 \(f(x)=f(y)\)，
否则 \(F\) 根本不是映射。类似地，把一个抽象对称性写进
物理计算以前，也须先给出它在状态空间上的表示。
这些看似形式的检查，正是后文处理波函数相位、算子定义域与坐标
变换时避免矛盾的第一道关口。
